\chapter{数列}

\section{数列的概念}

\begin{definition}[数列]
按照一定顺序排列的一列数叫做\textbf{数列}。数列中的每一个数叫做这个数列的\textbf{项}，第 $n$ 项记作 $a_n$。
\end{definition}

\begin{definition}[数列的通项公式]
如果数列 $\{a_n\}$ 的第 $n$ 项 $a_n$ 与序号 $n$ 之间的关系可以用一个公式表示，那么这个公式就叫做这个数列的\textbf{通项公式}。
\end{definition}

\begin{definition}[数列的前 $n$ 项和]
数列 $\{a_n\}$ 的前 $n$ 项和记作 $S_n$，即：
\[ S_n = a_1 + a_2 + \cdots + a_n = \sum_{k=1}^{n} a_k \]
\end{definition}

\begin{theorem}[通项与前 $n$ 项和的关系]
\[ a_n = \begin{cases} S_1, & n = 1 \\ S_n - S_{n-1}, & n \geq 2 \end{cases} \]
\end{theorem}

\section{等差数列}

\begin{definition}[等差数列]
如果一个数列从第二项起，每一项与它的前一项的差都等于同一个常数，这个数列就叫做\textbf{等差数列}。这个常数叫做等差数列的\textbf{公差}，通常用字母 $d$ 表示。
\end{definition}

\begin{theorem}[等差数列的通项公式]
设等差数列 $\{a_n\}$ 的首项为 $a_1$，公差为 $d$，则：
\[ a_n = a_1 + (n-1)d \]
\end{theorem}

\begin{theorem}[等差数列的前 $n$ 项和公式]
\[ S_n = \frac{n(a_1 + a_n)}{2} = na_1 + \frac{n(n-1)}{2}d \]
\end{theorem}

\begin{note}
记忆口诀：\textbf{首末相加乘项数，除以二}——等差数列求和公式的形象记忆。
\end{note}

\begin{property}[等差数列的性质]
\begin{enumerate}
    \item 若 $m + n = p + q$，则 $a_m + a_n = a_p + a_q$；
    \item 等差中项：$2a_n = a_{n-1} + a_{n+1}$（$n \geq 2$）；
    \item 数列 $\{a_n\}$ 是等差数列 $\Leftrightarrow$ $a_n = An + B$（$A, B$ 为常数）。
\end{enumerate}
\end{property}

\begin{example}
已知等差数列 $\{a_n\}$ 中，$a_3 = 5$，$a_7 = 13$，求 $a_{10}$ 和 $S_{10}$。
\end{example}

\textbf{解.}

\textbf{第一步：求公差}
由 $a_7 = a_3 + 4d$，得 $13 = 5 + 4d$，所以 $d = 2$.

\textbf{第二步：求首项}
$a_1 = a_3 - 2d = 5 - 4 = 1$.

\textbf{第三步：求 $a_{10}$}
$a_{10} = a_1 + 9d = 1 + 18 = 19$.

\textbf{第四步：求 $S_{10}$}
\[ S_{10} = \frac{10(a_1 + a_{10})}{2} = \frac{10(1 + 19)}{2} = 100 \]
\hfill $\square$

\section{等比数列}

\begin{definition}[等比数列]
如果一个数列从第二项起，每一项与它的前一项的比都等于同一个常数，这个数列就叫做\textbf{等比数列}。这个常数叫做等比数列的\textbf{公比}，通常用字母 $q$ 表示（$q \neq 0$）。
\end{definition}

\begin{theorem}[等比数列的通项公式]
设等比数列 $\{a_n\}$ 的首项为 $a_1$，公比为 $q$，则：
\[ a_n = a_1 \cdot q^{n-1} \]
\end{theorem}

\begin{theorem}[等比数列的前 $n$ 项和公式]
\[ S_n = \begin{cases} na_1, & q = 1 \\ \dfrac{a_1(1-q^n)}{1-q} = \dfrac{a_1 - a_nq}{1-q}, & q \neq 1 \end{cases} \]
\end{theorem}

\begin{note}
注意：使用等比数列求和公式时，必须先判断公比 $q$ 是否为 $1$！
\end{note}

\begin{property}[等比数列的性质]
\begin{enumerate}
    \item 若 $m + n = p + q$，则 $a_m \cdot a_n = a_p \cdot a_q$；
    \item 等比中项：$a_n^2 = a_{n-1} \cdot a_{n+1}$（$n \geq 2$，各项同号时）；
    \item 等比数列中不能有 $0$ 项。
\end{enumerate}
\end{property}

\begin{example}
已知等比数列 $\{a_n\}$ 中，$a_1 = 2$，$q = 3$，求 $a_5$ 和 $S_5$。
\end{example}

\textbf{解.}

\[ a_5 = a_1 \cdot q^4 = 2 \cdot 3^4 = 2 \cdot 81 = 162 \]

\[ S_5 = \frac{a_1(1-q^5)}{1-q} = \frac{2(1-243)}{1-3} = \frac{2(-242)}{-2} = 242 \]
\hfill $\square$

\section{数列求和方法}

\begin{theorem}[裂项相消法]
常见裂项公式：
\begin{enumerate}
    \item $\dfrac{1}{n(n+1)} = \dfrac{1}{n} - \dfrac{1}{n+1}$
    \item $\dfrac{1}{n(n+2)} = \dfrac{1}{2}\left(\dfrac{1}{n} - \dfrac{1}{n+2}\right)$
    \item $\dfrac{1}{(2n-1)(2n+1)} = \dfrac{1}{2}\left(\dfrac{1}{2n-1} - \dfrac{1}{2n+1}\right)$
\end{enumerate}
\end{theorem}

\begin{theorem}[错位相减法]
适用于等差数列与等比数列对应项相乘构成的数列求和。

设 $S_n = a_1b_1 + a_2b_2 + \cdots + a_nb_n$，其中 $\{a_n\}$ 为等差数列，$\{b_n\}$ 为等比数列。
\end{theorem}

\begin{example}
求数列 $\left\{\dfrac{1}{n(n+1)}\right\}$ 的前 $n$ 项和 $S_n$。
\end{example}

\textbf{解.}

\[ \frac{1}{n(n+1)} = \frac{1}{n} - \frac{1}{n+1} \]

所以：
\begin{align*}
S_n &= \left(1 - \frac{1}{2}\right) + \left(\frac{1}{2} - \frac{1}{3}\right) + \left(\frac{1}{3} - \frac{1}{4}\right) + \cdots + \left(\frac{1}{n} - \frac{1}{n+1}\right) \\
&= 1 - \frac{1}{n+1} = \frac{n}{n+1}
\end{align*}
\hfill $\square$
